\section{A conditional computation rather than a verbal certificate}
\label{sec:C09:conditional}
Inference-time reasoning changes the computation used to choose an answer or action while leaving some or all model weights fixed. A model can produce intermediate text, sample alternatives, invoke tools, or delegate subproblems. These mechanisms have different information and resource effects. Intermediate text extends the model's conditional state; an external tool can supply new evidence; independent sampling explores alternative continuations; a verifier changes which continuation is selected. Calling all of them thinking obscures the intervention. Their contribution should be identified through the distribution they change and the evidence they add.

Let $z$ denote an intermediate reasoning sequence and $a$ the final action. A generative model defines $p_\theta(z,a\mid C)$, and its induced action distribution marginalizes over $z$. Choosing one greedy sequence, sampling a sequence, voting across answers, and searching a tree produce different decision rules from the same underlying model. They need not rank actions identically. A system that reports only the final action without the generation and selection procedure does not fully specify its inference policy. The selected reasoning trace also need not characterize all alternative computations the model could have produced.

Chain-of-thought prompting supplies examples or instructions that elicit intermediate language before an answer, with reported improvements in the original study's evaluated reasoning tasks~\cite{wei2022cot}. This is an empirical prompting result, not a guarantee that every generated rationale is a faithful account of internal computation. A rationale can contain a correct derivation, a plausible explanation of a guessed answer, or an error later corrected. For an agent, the distinction becomes consequential when a trace is stored as procedural memory. A fluent but invalid derivation should not become trusted guidance merely because the final answer happened to be correct.

The strongest reasoning evidence is an independently checkable object attached to a precise claim: an executable calculation, a proof accepted by a specified checker, a test demonstrating a behavior, or a trace of actual tool observations. Even these objects have scope limits. A proof checker validates a formal statement under its axioms; a test covers selected inputs; a calculator evaluates the supplied expression rather than determining whether it models the user's problem. Reasoning quality therefore includes translating the task into the right formal object and interpreting its result, not only generating a long sequence of intermediate tokens.

